% !TEX TS-program = lualatex
%\documentclass[handout]{beamer}
\documentclass{beamer}
\input{../common/misomva}
\newcommand{\misoclasstinytitleslide}{Partially based on material by:  Ulrike von Luxburg, \\[-1em] Gary Miller, Doyle \& Schnell, Daniel Spielman}
\newcommand{\misoclassdate}{}
\usetikzlibrary{graphs}

\begin{document}
% Topic-specific title slide

\newcommand{\topictitle}{Laplacian and Random Walks}

\newcommand{\topicsubtitle}{Stationary Distribution}

\MakeTopicSlide

\begin{frame}
\frametitle{Normalized Laplacians}

\begin{block}{}
$\bL_{sym}$ and $\bL_{rw}$ are PSD with non-negative real eigenvalues 
\vspace{-0.25cm}
$$0 = \lambda_1 \leq \lambda_2 \leq \lambda_3 \leq \dots \leq \lambda_\nodes.$$
\vspace{-0.5cm} 
\end{block}


\begin{block}{}
$(\lambda,\bu)$ is an eigenpair for $\bL_{rw}$ iff $(\lambda,\bu)$
solve the generalized eigenproblem $\bL \bu = \lambda \bD \bu$.
\end{block}\pause

\begin{block}{}
$(0,\bOne_\nodes)$ is an eigenpair for $\bL_{rw}$.
\end{block}\pause

\begin{block}{}
$(0,\bD^{1/2}\bOne_\nodes)$ is an eigenpair for $\bL_{sym}$.
\end{block}\pause

\begin{block}{}
Multiplicity of eigenvalue 0 of $\bL_{rw}$ or $\bL_{sym}$ equals
to the number of connected components.
\end{block}\pause
\begin{flushright}
\notecol{Proof: As for $\bL$.}\end{flushright}

\end{frame}
%%%%%%%%%%%%%%%%%%%%%%%%%%%%%%%%%%%%%%%%%%%%%%%%%%%%%%%%%%%%%
\begin{frame}
\frametitle{Laplacian and Random Walks on Undirected Graphs}
\begin{itemize}[<+-| alert@+>]\addtolength{\itemsep}{0.5\baselineskip}
\item stochastic process: vertex-to-vertex jumping
\item transition probability $v_i \to v_j$ is $p_{ij} = w_{ij}/d_i$ 
\begin{itemize}
\item $d_i \eqdef \sum_j w_{ij} $
\end{itemize}
\item transition matrix $\bP = (p_{ij})_{ij} = \bD^{-1}\bW$ (notice $\bL_{rw} = \bI - \bP$)
\item if $G$ is connected and non-bipartite $\to$ unique 
\textbf{stationary distribution} 
$\pi = (\pi_1,\pi_2,\pi_3,\dots,\pi_\nodes)$ where $\pi_i = d_i/{\rm vol}(V)$
\begin{itemize}
\item ${\rm vol}(G) = {\rm vol}(V) = {\rm vol}(\bW) \eqdef \sum_i d_i = \sum_{i,j} w_{ij}$
\end{itemize}
\item $ \pi = \frac{{\mathbf 1\transpose}\bW}{\mathrm{vol}(\bW)}$ verifies $\pi \bP = \pi$ as:
\end{itemize}
\pause
\begin{align*}
\pi  \bP  &=  \frac{{\mathbf 1\transpose}\bW\bP}{\mathrm{vol}(\bW)}= \frac{{\mathbf
1\transpose}\bD\bP}{\mathrm{vol}(\bW)} = \frac{{\mathbf 1\transpose}\bD\bD^{-1}\bW}{\mathrm{vol}(\bW)} =
\frac{{\mathbf 1\transpose}\bW}{\mathrm{vol}(\bW)}= \pi
%s  &=  {\mathbf 1}D/\mathrm{vol}(D) = {\mathbf 1}W/\mathrm{vol}(W)
\end{align*} \pause
\rightquestionbox{What's the difference from the \pagerank\!\!\!\textsuperscript{TM}?}
\end{frame}


\MakeLastSlide
\end{document}
